\section{Background: Narrowing the Label Space with Conformal Prediction}
\label{sec:background}

\paragraph{Split conformal prediction, in plain terms.}
Take any classifier that scores each class.
Hold out a calibration split it has not seen and record, for each calibration example, how low the classifier scored the true label.
For a new input, include in the candidate set every class whose score is no worse than what the true label received on all but a fraction $\alpha$ of the calibration examples.
If the new input is exchangeable with the calibration split, the true label is in the set with probability at least $1-\alpha$, for any classifier and any sample size \citep{vovk2005algorithmic,papadopoulos2002inductive}.
A weak classifier does not lower the coverage, it makes the sets larger.

\paragraph{Class-conditional calibration.}
The guarantee above is marginal over classes.
On a long-tailed dataset the set can cover 95\% of inputs overall while missing most inputs of a rare class, because the rare class contributes few calibration examples to the shared threshold.
Class-conditional conformal prediction computes one threshold per class from the $n_c$ calibration examples of class $c$ only, which guarantees $P(y \in \mathcal{C}(x) \mid y = c) \geq 1-\alpha$ for every class, in expectation over the draw of the calibration split \citep{vovk2012conditional}.
The coverage realised on a finite test sample fluctuates around this level, more so when $n_c$ is small, and with fewer than $1/\alpha - 1$ calibration examples (19 at $\alpha = 0.05$) the class is kept in almost every set.
Top-$k$ selection and probability thresholds give no such guarantee, the first because it ignores the classifier's confidence, the second because it trusts a confidence that is rarely calibrated.

\paragraph{The CICLe pipeline.}
CICLe \citep{randl2024cicle} trains a base classifier (logistic regression on sentence embeddings) on the labelled pool and calibrates it class-conditionally.
For a test input it yields a candidate set $\mathcal{C}(x)$.
If $|\mathcal{C}(x)| = 1$ that label is returned without calling the LLM, and if the set is empty the classifier's top label is.
Otherwise the LLM is prompted with the task description, the candidate labels and labelled examples retrieved from the candidate classes only.

\paragraph{Narrowing rules compared.}\label{sec:rules}
To separate narrowing as such from conformal calibration, we build candidate sets from the same classifier with five rules.
\emph{CP} is the class-conditional conformal set at miscoverage $\alpha$.
\emph{Marginal CP} shares one threshold across classes.
\emph{Top-$m$} keeps the $m$ most probable classes, where $m$ is the mean CP set size on that dataset and seed, so the budget matches but the size does not adapt to the input.
\emph{Probability mass} keeps the most probable classes until their predicted probabilities sum to $1-\alpha$, which adapts but trusts the uncalibrated confidence.
\emph{Oracle} contains the gold label plus $m-1$ random classes and bounds any rule at that size.
All rules share the prompt, the retrieval and the singleton and empty-set behaviour.

\paragraph{Fixed and Per-Class retrieval.}
\emph{Fixed} retrieval places the $k$ training examples most similar to the input (cosine similarity in the embedding space) among those whose label is in the candidate set.
\emph{Per-Class} retrieval places the $k$ most similar examples of each candidate class, so the prompt holds $k \times |\mathcal{C}(x)|$ examples and shrinks with the set.
Few-shot prompting is the case in which the candidate set is the full label set.
Narrowing also makes listing labels affordable, since few-shot prompting must list every class and CICLe only the candidate set.
